\chapter{\IfLanguageName{dutch}{Requirements en Ontwerp}{Requirements and Design}}%
\label{ch:requirements-ontwerp}

Uit de voorgaande Salesforce analyse lijkt de native In-App Guidance uitstekend geschikt om een gebruiker te begeleiden bij allerlei soorten processen. Voor de initiële setup van een complexe applicatie zoals 24flow moet deze guidance ook duidelijk georganiseerd worden. Daarom wordt er voor het proof of concept gewerkt met een op maat gemaakt 'Onboarding Dashboard' waar de gebruiker op landt bij het starten van de applicatie. Dit dashboard fungeert als een centrale hub die de gebruiker een duidelijk stappenplan en houvast biedt. Vanuit dit overzicht kan de gebruiker dan stap per stap de nodige configuraties opzetten met behulp van de in-app guidance zonder het overzicht te verliezen en met een duidelijke houvast.

\section{Requirements Analyse}%
\label{sec:Requirements analyse}
Voor de technische uitwerking van deze oplossing is het cruciaal om de vereisten van het systeem helder in kaart te brengen. In deze sectie worden de kerncomponenten van 24flow, de doelgroep, de niet-functionele requirements en de specifieke functionele requirements gedefinieerd en geprioritiseerd.

\subsection{Kerncomponenten van 24flow}
\label{subsec:kerncomponenten_24flow}
Om de complexiteit van de onboarding te begrijpen, is inzicht in de modulaire opbouw van 24flow noodzakelijk. Voor dit onderzoek wordt er enkel gefocust op één van de verschillende modules die 24flow aanbiedt, namelijk 24factory, waar er een digitale fabriek gebruikt wordt die de effectieve werkvloer beheert van een klant. Dit platform bestaat uit verschillende interactieve componenten die nauw met elkaar verbonden zijn via Salesforce-data:

\begin{itemize}
	\item \textbf{Order Pool:} Een centraal overzicht waar alle productieorders binnenkomen. Deze component fungeert als de ``ingang'' van de fabriek. De effectiviteit van de Order Pool hangt af van de juiste filtering en weergave van velden (zoals leverdatum en producttype). 
	\item \textbf{Team Cockpit:} Een visueel bord waarop operatoren de status van Productie Orders (bijv. \textit{In Buffer, In Process, Expected}) kunnen opvolgen. Dit is het hart van de werkvloer planning, en hiermee kan de hoeveelheid work in progress op de vloer georganiseerd worden.
	\item \textbf{Production Orders, Teams \& Operations:} Dit zijn enkele van de belangrijkste achterliggende datastructuren. Een productieorder bestaat uit meerdere bewerkingen (operations) die door verschillende teams uitgevoerd moeten worden.
\end{itemize}

\subsubsection*{De afhankelijkheid van data}
Een fundamenteel kenmerk van deze componenten is dat zij \textbf{visueel georiënteerd} zijn. Een Order Pool zonder records is een lege tabel, en een Team Cockpit zonder actieve bewerkingen toont enkel lege kolommen.
\\

Daarom is er de noodzaak om een gebruiker over realistische data te laten beschikken tijdens de onboarding. Ofwel zou er demo-omgeving gebruikt kunnen worden of een klant zou zijn eigen data kunnen uploaden en gebruiken of een combinatie van beide indien mogelijk. Dit zou ervoor zorgen dat de gebruiker direct met een gevulde interface kan werken, wat essentieel is voor een succesvolle configuratie.

\subsection{Doelgroep}
\label{subsec:doelgroepanalyse}
Voor de uitwerking van een effectieve onboarding-oplossing is het cruciaal om te bepalen voor wie het systeem wordt ontworpen. Binnen het 24flow-ecosysteem zijn er twee specifieke gebruikersrollen geïdentificeerd. Hoewel de uiteindelijke applicatie voor beide rollen hetzelfde is, verschillen hun doelen en vereisten:

\begin{itemize}
	\item \textbf{De Salesforce-beheerder van de klant:} Dit is een beheerder bij een (potentieel) nieuw klantbedrijf die 24flow wil evalueren of implementeren. Deze gebruiker kent de applicatie nog niet en wil het systeem zelfstandig leren gebruiken en testen. Voor deze groep is er een sterke behoefte aan stapsgewijze, interactieve begeleiding en direct toegankelijke en duidelijke documentatie. Daarnaast is het voor hen belangrijk om over realistische testdata te beschikken, zodat zij de applicatie direct visueel kunnen evalueren zonder eerst uren aan handmatige data-invoer te spenderen. Het uiteindelijke doel van deze gebruiker is 24flow in werking zien en indien ze nog geen klant zijn evalueren of ze wel of niet klant willen worden.
	
	\item \textbf{De Solution Engineer (24flow-expert):} Dit is een interne medewerker van 24flow. Hoewel het onboarding-systeem primair bedoeld is voor (nieuwe) klanten, bieden specifieke onderdelen ervan ook een grote meerwaarde voor deze interne experts. Zij kennen de applicatie al door en door en hebben geen behoefte aan stapsgewijze handleidingen. Echter, de functionaliteiten die nodig zijn om een omgeving geautomatiseerd op te zetten en van testdata te voorzien, kunnen door hen uitstekend worden ingezet om in zeer korte tijd bijvoorbeeld een demo-omgeving klaar te zetten.
\end{itemize}

\subsection{Centralisatie}
\label{subsec:centralisatie_dashboard}
Een van de grootste uitdagingen bij een standaard Salesforce-implementatie is de verspreiding van de instellingen. Omdat de benodigde configuraties verspreid liggen over talloze menu's, objecten en Setup-pagina's, ontstaat er al snel een gefragmenteerde gebruikerservaring. Om te voorkomen dat een nieuwe gebruiker het overzicht verliest, is er een duidelijke noodzaak voor centralisatie. 
\\

Het te ontwikkelen systeem moet daarom een centraal vertrekpunt bevatten dat voor structuur bij de eindgebruiker kan zorgen. In plaats van de gebruiker te verplichten zelf op zoek te gaan naar alle componenten, moeten alle noodzakelijke onboarding-taken op deze centrale plaats worden samengebracht. Dit vermijdt onnodig zoekwerk en zorgt ervoor dat een gebruiker het overzicht steeds bewaart. Hier kan dan ook een link naar de documentatie zijn en dit kan eveneens alle extra tools bevatten die nodig zouden kunnen zijn voor een goede onboarding.

\subsection{Niet-Functionele Requirements}
De oplossing moet voldoen aan de volgende technische randvoorwaarden binnen het Salesforce-platform:

\begin{itemize}
	\item \textbf{Interface:} De gebruikersinterface moet gebouwd worden conform het Salesforce Lightning Design System (SLDS).
	\item \textbf{Performantie:} Alle backend-processen (zoals Apex data-generatie) moeten binnen de vastgestelde \textit{governor limits} van Salesforce opereren.
	\item \textbf{Beheer:} De configuratie van de in-app guidance (prompts en walkthroughs) moet aanpasbaar zijn via de standaard Salesforce set-up, zonder dat er codewijzigingen nodig zijn.
\end{itemize}

\subsection{Prioritering van de requirements}
Om de scope van het proof of concept onder controle te houden en te garanderen dat de meest kritische onderdelen als eerste worden gebouwd, zijn de verzamelde requirements geprioriteerd volgens de MoSCoW-methode.

\subsubsection*{Must haves}
De "Must haves" vormen de minimale set aan functionaliteiten die aanwezig moeten zijn om een werkbare oplossing te bieden die beantwoordt aan de onderzoeksvraag:

\begin{itemize}
	\item \textbf{Zelfstandige Systeemconfiguratie:} Een volledige interactieve onboarding die de gebruiker stap-voor-stap door de initiële opzet van 24flow loodst voor zijn belangrijkste componenten. Dit stelt de gebruiker in staat om de applicatie volledig operationeel te krijgen zonder externe hulp.
	\item \textbf{Centraal Dashboard:} Een overzichtelijk startpunt waar de gebruiker de voortgang van zijn configuratie kan monitoren via een visuele voortgangsbalk en status-indicatoren. Dit is eveneens de landingspagina van elke 24flow-applicatie.
	\item \textbf{Directe Navigatie en Activatie:} Een koppeling waarbij de gebruiker met één druk op de knop direct naar de juiste pagina wordt gestuurd en de bijbehorende  walkthrough automatisch wordt gestart.
	\item \textbf{Automatische Voortgangsdetectie:} Het systeem moet zelfstandig herkennen wanneer een taak succesvol is afgerond en dit direct verwerken in het dashboard. Het moet per onderdeel bijhouden of en in hoeverre de configuratie is voltooid.
	\item \textbf{Demo Data Generator:} De mogelijkheid om de omgeving te vullen met actuele testdata. Dit is essentieel voor een realistisch gebruik van de applicatie, zodat de gebruiker de onboarding doorloopt in een "gevulde" interface.
	\item \textbf{Auto-Configuratie:} Het systeem moet waar het kan zoveel mogelijk automatisch de benodigde Lightning Pages en basisconfiguraties kunnen genereren voor de kerncomponenten. Dit minimaliseert de tijd voor de gebruiker om tot een werkende omgeving te komen.
	\item \textbf{Hiërarchische Taakstructuur:} De mogelijkheid om grote, complexe processen op te delen in kleinere subtaken, zodat de gebruiker stap-voor-stap door de configuratie wordt geleid zonder het overzicht te verliezen.
	\item \textbf{Navigatie tussen objecten:} De mogelijkheid voor walkthroughs om de gebruiker naadloos over verschillende tabbladen en records heen te leiden, aangezien deze zelden beperkt blijven tot één enkel scherm.
	\item \textbf{Documentatie:} Een directe en permanente link naar de officiële documentatie op het dashboard. Dit dient als hulpmiddel en naslagwerk voor gebruikers die specifieke, of meer diepgaande informatie zoeken over bepaalde componenten.
\end{itemize}

\subsubsection*{Should haves}
De "Should haves" zijn belangrijke functionele aanvullingen die de kwaliteit van de onboarding verhogen en essentieel zijn voor een volledige en professionele gebruikerservaring:

\begin{itemize}
	\item \textbf{Video-integratie:} Het toevoegen van korte video-tutorials in docked prompts voor de meest complexe processen. Dit verhoogt de effectiviteit van de onboarding, maar kan in eerste instantie worden vervangen door tekst en afbeeldingen.
	\item \textbf{Data Reset-functie:} Een veilige methode om testdata te wissen, zodat een gebruiker altijd met een propere omgeving kan werken.
	\item \textbf{Uploaden Eigen Data:} Een methode om eigen data te uploaden die gebruikt kan worden voor de onboarding. Dit kan in eerste instantie een placeholder zijn die later kan worden geïmplementeerd.
	\item \textbf{Taakbevestiging:} Voor stappen die niet automatisch door het systeem gemeten kunnen worden, moet de gebruiker zelf een taak als voltooid kunnen markeren.
	\item \textbf{Custom Branding:} Het aanpassen van de kleuren van de prompts aan de huisstijl van 24flow. Dit draagt bij aan de uniformiteit en herkenbaarheid binnen het platform.
	\item \textbf{Documentatie Linken} Getoonde documentatie moet specifiek gekoppeld zijn aan de respectievelijke component. Een gebruiker mag niet op zoek moeten gaan in alle documentatie voor een specifiek onderdeel.
\end{itemize}

\subsubsection*{Could haves}
Deze elementen worden als "extra" beschouwd en worden enkel geïmplementeerd als er na de kritieke onderdelen nog tijd over is:

\begin{itemize}
	\item \textbf{Hervatten:} De optie voor gebruikers om een reeds voltooide walkthrough op elk moment opnieuw te starten en verder te gaan op het punt waar ze gebleven zijn.
	\item \textbf{Gebruikersstatistieken (Analytics):} Het opzetten van rapportage om te meten hoe vaak walkthroughs worden voltooid en waar gebruikers vastlopen. Dit is nodig voor de evaluatie op lange termijn, maar niet voor de technische werking.
	\item \textbf{Gebruikersfeedback:} De mogelijkheid voor gebruikers om direct binnen een prompt aan te geven of de instructie nuttig was via een simpele duimpje-omhoog/omlaag interactie.
	\item \textbf{Gamificatie-elementen:} Het toevoegen van visuele elementen, zoals een vuurwerk animatie of een digitale "badge" bij het voltooien van een bepaalde configuratie, om de gebruikerservaring te verbeteren.
	\item \textbf{Zoekfunctionaliteit:} Een zoekbalk op het dashboard waarmee gebruikers direct naar specifieke instellingen of documentatie-onderwerpen kunnen springen.
	\item \textbf{Directe Support-link:} Een knop op het dashboard die direct een support-ticket opent bij een 24flow-expert voor vragen die niet door de onboarding worden beantwoord.
\end{itemize}

\subsubsection*{Won't haves}
Deze onderdelen zijn expliciet buiten de scope van dit project gelaten, vaak vanwege technische beperkingen van het Salesforce-platform:

\begin{itemize}
	\item \textbf{Mobiele Ondersteuning:} Ondersteuning voor de Salesforce Mobile App, aangezien in-app guidance momenteel technisch enkel beschikbaar is voor de desktopomgeving.
	\item \textbf{Chatbot:} AI Chatbots: Een interactieve bot die vragen beantwoordt; we focussen op gestuurde navigatie, de gebruiker moet in staat zijn om zelf wegwijs te worden op het platform.
\end{itemize}

\section{Ontwerp proof of concept}
\label{sec:Ontwerp}

\subsection{Use cases}

\subsubsection{Use Case: Initiële configuratie van de Order Pool}
Deze use case beschrijft hoe een nieuwe gebruiker zelfstandig een essentieel onderdeel van de applicatie inricht via de onboarding oplossing.
\\

Actor: Nieuwe 24flow-gebruiker (Systeembeheerder/Admin).
\\

Doel: Het succesvol en zelfstandig configureren van de Order Pool-instellingen.

\begin{enumerate}
	\item \textbf{Startpunt:} De gebruiker landt op de centrale landingspagina van de applicatie. Het \textbf{Onboarding Dashboard} verschijnt direct in beeld.
	\item \textbf{Status:} De gebruiker ziet in de voortgangsbalk dat de setup op 0 procent staat. In de hiërarchische takenlijst wordt de taak "Order Pool instellen" als openstaand gemarkeerd.
	\item \textbf{Activatie:} De gebruiker klikt op de actieknop naast de taak. Het systeem navigeert de gebruiker direct naar de juiste instellingenpagina van de Order Pool.
	\item \textbf{Begeleiding:} Bij het laden van de pagina start automatisch de Salesforce \textbf{Walkthrough}. Een targeted prompt wijst naar het eerste configuratieveld en legt uit welke informatie hier benodigd is.
	\item \textbf{Interactie:} De gebruiker volgt de stappen van de walkthrough, waarbij hij door verschillende tabbladen wordt geleid om filters en kolommen te definiëren.
	\item \textbf{Detectie:} Nadat de gebruiker op 'Opslaan' heeft geklikt, herkent de \textbf{Automatische Voortgangsdetectie} op de achtergrond dat de configuratie-records zijn aangemaakt.
	\item \textbf{Voltooiing:} De gebruiker keert terug naar het dashboard. De taak "Order Pool instellen" is nu voorzien van een groen vinkje en de algehele voortgangsbalk is bijgewerkt naar een hoger percentage.
\end{enumerate}

\subsubsection{Use Case: Beheer van testdata via het dashboard}
Deze use case illustreert hoe de functionaliteiten in de zijbalk van het dashboard worden ingezet om de applicatie klaar te maken voor gebruik door testdata te genereren.
\\

Actor: Nieuwe 24flow-gebruiker (Systeembeheerder/Admin).
\\

Doel: Het succesvol vullen van de applicatie met realistische testdata, zodat de functionaliteiten optimaal kunnen worden geconfigureerd en gebruikt.

\begin{enumerate}
	\item \textbf{Startpunt:} De gebruiker bevindt zich op het Onboarding Dashboard en merkt op dat de applicatie nog leeg is en geen data bevat.
	\item \textbf{Actie:} In het "Tools" paneel aan de rechterkant klikt de gebruiker op de knop "Generate Demo Data".
	\item \textbf{Verwerking:} Het systeem toont tijdelijk een laad-indicator. Op de achtergrond wordt een Apex-proces aangeroepen dat automatisch een set gestandaardiseerde testrecords aanmaakt.
	\item \textbf{Bevestiging:} Zodra het proces is afgerond, toont het dashboard een succesmelding (toast notificatie) aan de gebruiker.
	\item \textbf{Resultaat:} De applicatie is nu gevuld met data, waardoor de gebruiker bij walkthroughs direct deze data kan gebruiken om de nodige configuraties te doen.
\end{enumerate}

\clearpage
\subsection{Wireframes}
Om de requirements visueel te vertalen, zijn er enkele wireframes uitgewerkt. Dit ontwerp dient als blauwdruk voor de uiteindelijke Lightning Web Component en de onboarding implementatie.

\subsubsection{Wireframe: Onboarding Dashboard}
Zoals te zien is in Figuur \ref{fig:wireframe_dashboard}, fungeert het dashboard als de landingspagina en het vertrekpunt voor een nieuwe gebruiker. Aan de linkerkant bevindt zich de hiërarchische taakstructuur met dynamische status-indicatoren en een voortgangsbalk. Aan de rechterkant zijn enkele tools om bijvoorbeeld data te genereren en wissen en documentatie te raadplegen.
\begin{figure}[htbp]
	\centering
	\includegraphics[width=0.95\textwidth]{WireframeDashboardOnboarding}
	\caption{Wireframe van het Onboarding Dashboard inclusief voortgangsbalk en zijbalk met extra tools.}
	\label{fig:wireframe_dashboard}
\end{figure}

\clearpage
\subsubsection{Wireframe: Walkthrough Stap}
Wanneer de gebruiker in het dashboard een taak start, navigeert het systeem direct naar de desbetreffende configuratiepagina. Figuur \ref{fig:wireframe_walkthrough} toont deze stap. Hier neemt de in-app guidance het proces over in de vorm van een walkthrough. Een \textit{Targeted Prompt} wijst de gebruiker expliciet op het in te vullen veld en biedt direct een aanbeveling ("Action Link") voor meer info.
\\

Tegelijkertijd biedt een \textit{Docked Prompt}, onderaan het scherm, extra ondersteuning via een instructievideo en een documentatielink. Dit ontwerp zorgt ervoor dat de gebruiker complexe acties stapsgewijs en zelfstandig kan uitvoeren en ook steeds een directe link heeft naar extra documentatie indien nodig.

\begin{figure}[htbp]
	\centering
	\includegraphics[width=0.95\textwidth]{WalkthroughStep}
	\caption{Wireframe van een walkthrough stap voor het configureren van de Order Pool.}
	\label{fig:wireframe_walkthrough}
\end{figure}